\documentclass[11pt]{article}

\usepackage[a4paper,margin=27mm]{geometry}
\usepackage{amsmath,amssymb,amsthm,mathtools}
\usepackage[T1]{fontenc}
\usepackage{lmodern}
\usepackage{microtype}
\usepackage{array,tabularx}
\usepackage{needspace,float,aliascnt}
\usepackage{booktabs}
\usepackage{enumitem}
\usepackage{xcolor}
\usepackage[hidelinks]{hyperref}
\hypersetup{pdftitle={Exact-Value Complexity of Nonsmooth Strongly Convex Optimization},
  pdfauthor={Colombo Team},pdfsubject={Uniform oracle complexity bounds and boundary geometry}}
\setlength{\emergencystretch}{1.5em}
\setlength{\parskip}{2pt}
\setlist{itemsep=3pt,topsep=5pt}
\newcolumntype{Y}{>{\raggedright\arraybackslash}X}
\newcolumntype{P}[1]{>{\raggedright\arraybackslash}p{#1}}
\newcommand{\logfac}{\mathcal L_\eps}
\newcommand{\profile}[1]{\mathcal P_d(#1)}
\usepackage[nameinlink,capitalize]{cleveref}

\newcommand{\R}{\mathbb{R}}
\newcommand{\E}{\mathbb{E}}
\newcommand{\Pp}{\mathbb{P}}
\newcommand{\B}{\mathbb{B}}
\newcommand{\Sph}{\mathbb{S}}
\newcommand{\cF}{\mathcal{F}}
\newcommand{\cA}{\mathcal{A}}
\newcommand{\cD}{\mathcal{D}}
\newcommand{\eps}{\varepsilon}
\newcommand{\inner}[2]{\langle #1,#2\rangle}
\newcommand{\norm}[1]{\lVert #1\rVert_2}
\newcommand{\ind}{\mathbf{1}}
\newcommand{\polylog}{\operatorname{polylog}}
\newcommand{\Lip}{\operatorname{Lip}}
\newcommand{\conv}{\operatorname{conv}}
\newcommand{\val}{\operatorname{val}}
\newcommand{\err}{\operatorname{err}}
\newcommand{\diag}{\operatorname{diag}}
\newcommand{\argmin}{\operatorname*{arg\,min}}
\newcommand{\argmax}{\operatorname*{arg\,max}}

\theoremstyle{plain}
\newtheorem{theorem}{Theorem}[section]
\newaliascnt{proposition}{theorem}
\newtheorem{proposition}[proposition]{Proposition}
\aliascntresetthe{proposition}
\newaliascnt{lemma}{theorem}
\newtheorem{lemma}[lemma]{Lemma}
\aliascntresetthe{lemma}
\newaliascnt{corollary}{theorem}
\newtheorem{corollary}[corollary]{Corollary}
\aliascntresetthe{corollary}
\theoremstyle{definition}
\newaliascnt{definition}{theorem}
\newtheorem{definition}[definition]{Definition}
\aliascntresetthe{definition}
\newaliascnt{remark}{theorem}
\newtheorem{remark}[remark]{Remark}
\aliascntresetthe{remark}
\crefname{proposition}{Proposition}{Propositions}
\crefname{lemma}{Lemma}{Lemmas}
\crefname{corollary}{Corollary}{Corollaries}
\crefname{definition}{Definition}{Definitions}
\crefname{remark}{Remark}{Remarks}

\title{Exact-Value Complexity of Nonsmooth\\
Strongly Convex Optimization:\\
Uniform Bounds and Boundary Geometry}
\author{Colombo Team\\[0.3em]
\small\href{mailto:colombo.2026@outlook.com}{\texttt{colombo.2026@outlook.com}}}
\date{}

\begin{document}
\maketitle

\begin{abstract}
We study the number of exact scalar function evaluations needed to minimize
a possibly nonsmooth, $\mu$-strongly convex, $L$-Lipschitz function on a
Euclidean ball of radius $R$. The relevant scale is the Lipschitz slack
$\Delta=L-\mu R$. We prove that the class is nonempty exactly when
$\Delta\ge0$ and collapses, at $\Delta=0$, to a known quadratic plus an
unknown constant. For $\Delta>0$, let $S_\Delta$ be the exact minimax
expected error without queries and set
\[
 m_\eps=\min\left\{d,\frac{\Delta^2R^2}{\eps^2},
                         \frac{\Delta^2}{\mu\eps}\right\}.
\]
For every dimension $d\ge1$ and $0<\eps\le c_0S_\Delta$, both expected-error
and constant-success query complexities lie between
$c d m_\eps/\log(e m_\eps)$ and
$C d m_\eps\log(d+1)\log(128S_\Delta/\eps)$.
The upper bound follows by localizing the unknown convex residual before
applying two-point estimation or value-oracle convex optimization.
The lower bound transfers a fixed-prior support-function obstruction
through quadratic regularization. A spherical family and a separate
one-dimensional packing complete the finite-dimensional cases. The large-dimensional
lower profile agrees with the Euclidean specialization of prior
strong-convexity reductions. Our contribution is a uniform geometric
and two-sided formulation, including the degenerate boundary, localized
logarithmic factors, and a direct regularization proof. The bounds
identify the polynomial dependence on dimension and accuracy; the
remaining logarithmic factors are not claimed to be sharp.
\end{abstract}

\paragraph{Keywords.}
zeroth-order optimization; exact function values; strong convexity;
nonsmooth optimization; oracle complexity; randomized algorithms.

\begin{samepage}
\noindent\textbf{AI Usage.}
Nearly the entire research pipeline for this paper was carried out by
Colombo, powered by GPT-5.6 Sol. 
\par
\end{samepage}

\section{Introduction}\label{sec:intro}

How many scalar evaluations are needed to minimize a strongly convex
function when derivatives are unavailable? In this paper, a query at
$x$ returns the exact real number $F(x)$ of one fixed objective. The
algorithm may be randomized and fully adaptive, and its internal
computation is unrestricted. The function may be nonsmooth. We count
individual evaluations and require every query and the final output
to lie in the prescribed Euclidean ball.

For a convex $L$-Lipschitz objective on a ball of radius $R$ and
$0<\eps\le LR$, two classical upper bounds are complementary.
The symmetric two-point method of
Shamir~\cite[Corollary~2]{Shamir2017}, specialized to a fixed objective,
uses $O(dL^2R^2/\eps^2)$ evaluations. The deterministic value-only method
of Protasov~\cite{Protasov1996} uses
$O(d^2\log(d+1)\log(4LR/\eps))$ evaluations. Lee, Sidford, and
Vempala~\cite{LeeSidfordVempala2018} give a broader evaluation-and-membership
oracle framework with a quadratic dimension dependence up to logarithms.
The precise logarithmic factors used here come from Protasov's bound.

Strong convexity suggests replacing the inverse-square accuracy
dependence by an inverse-linear one. On a fixed ball, however, the
parameters $L$, $\mu$, and $R$ cannot be treated independently. The
quadratic $\mu\|x\|_2^2/2$ already has radial slope $\mu R$ at the
boundary. The residual variation is therefore governed by
\begin{equation}\label{eq:slack-intro}
 \Delta:=L-\mu R.
\end{equation}
This distinction matters when $L$ approaches $\mu R$: the admissible
class then becomes progressively easier, and at equality every function
has the same minimizer, the origin. A uniform complexity statement must
capture this collapse.

\subsection{Results and proof ideas}

Our main result identifies, up to logarithmic factors, the polynomial
complexity profile
\begin{equation}\label{eq:profile-intro}
 d\min\left\{d,\frac{\Delta^2R^2}{\eps^2},
                  \frac{\Delta^2}{\mu\eps}\right\}.
\end{equation}
The statement holds in every finite dimension and at accuracies below
a universal constant times the error attainable without any queries.
Three ingredients make the dependence on $\Delta$ uniform.

\begin{enumerate}[label=(\roman*),leftmargin=2.2em]
\item \textbf{Geometry of the Lipschitz constraint.}
Writing $F(x)=h(x)+\mu\|x\|_2^2/2$, we show that each radial restriction
of $h$ is $\Delta$-Lipschitz, that the minimizer lies within distance
$\min\{R,\Delta/\mu\}$ of the origin, and that $h$ has an explicit
local Lipschitz bound. These facts identify the exact zero-query
expected-risk scale $S_\Delta$ and the rigid boundary $\Delta=0$.

\item \textbf{Localized algorithms.}
The geometry provides a known ball containing the minimizer on which
both the residual and the full objective have Lipschitz constants of
order $\Delta$. Two-point estimation gives the two accuracy-dependent
branches in \eqref{eq:profile-intro}. Applying Protasov's algorithm on
the same ball gives the dimension branch with a logarithm depending on
$S_\Delta/\eps$, rather than on the potentially much larger $LR/\eps$.

\item \textbf{A direct lower-bound transfer and finite-dimensional endpoints.}
We use the fixed support-function prior of Zhang et
al.~\cite[Theorem~5]{ZhangEtAl2026}. Subtracting a known quadratic
preserves its exact scalar transcript. A one-dimensional variational
profile then converts a poor output direction into objective error.
A hidden-direction quadratic family supplies the cases in which the
support-function calibration does not apply; a separate three-function
packing treats the one-dimensional constant-success endpoint.
\end{enumerate}

\subsection{Relation to previous lower bounds}\label{sec:related}

Zhang et al.~\cite{ZhangEtAl2026} establish the randomized convex
profile $d\min\{d,(LR/\eps)^2\}$ up to logarithms using a fixed random
support function. Kerger~\cite[Theorem~4]{Kerger2026} gives a
deterministic lower bound of order $d^2/\log(d+1)$ at accuracy
$cLR/\sqrt d$ and below. The fixed-prior result, rather than a
transcript-dependent resisting oracle, is the input needed here.

The closest strong-convexity comparison is Dorn et
al.~\cite[Theorem~K.2]{DornEtAl2026}. Their Euclidean specialization
already gives the large-dimensional lower profile
\eqref{eq:profile-intro} with the same slack $L-\mu R$. In their
notation, set $p=2$, $Q=I_d$, $m=d$, $r=1$, $a=1$, and $M=L$.
With
\[
 \sigma_\eps:=\min\left\{d,\frac49\min\left\{
       \frac{\Delta^2R^2}{\eps^2},\frac{\Delta^2}{\mu\eps}\right\}\right\},
\]
their bound is $c_Zd\sigma_\eps/\log(e\sigma_\eps)$, under
\begin{equation}\label{eq:dorn-window}
 d\ge m_Z^\circ,
 \qquad
 \frac{3\eps}{2}\le\eps_Z\Delta
                \min\left\{R,\sqrt{\frac{\eps}{\mu}}\right\}.
\end{equation}
Their stated criterion is success probability at least $2/3$.
The accuracy restriction is equivalent to
$\eps\le(2\eps_Z/3)\Delta R$ and
$\eps\le(2\eps_Z/3)^2\Delta^2/\mu$.

Table~\ref{tab:comparison} aligns these lower bounds by class, risk,
and parameter range. We do not claim a new large-dimensional polynomial
lower profile. We provide a uniform two-sided theorem, the associated
boundary geometry and localized upper bound, and an alternative
homogeneous-regularization proof. The scalar-information obstruction
continues to come from Zhang et al.

\begin{table}[H]
\centering
\small
\setlength{\tabcolsep}{4pt}
\renewcommand{\arraystretch}{1.22}
\caption{Comparison of exact-value lower bounds. Every row counts
individual feasible scalar evaluations. Here
$m_{\rm cvx}=\min\{d,(LR/\eps)^2\}$,
$\mathcal P_d(s)=ds/\log(es)$, and $m_\eps$ is defined in
\eqref{eq:m-eps}. Constants and dimension thresholds are universal.
Only the last row includes a matching upper bound in this table.}
\label{tab:comparison}
\begin{tabularx}{\textwidth}{@{}P{0.19\textwidth}P{0.22\textwidth}P{0.22\textwidth}Y@{}}
\toprule
Work & Class and criterion & Lower bound & Dimension and accuracy \\
\midrule
Kerger~\cite{Kerger2026}
& Convex; deterministic
& $\Omega(d^2/\log(d+1))$
& $d\ge d_0$, $\eps\le cLR/\sqrt d$ \\
Zhang et al.~\cite{ZhangEtAl2026}
& Convex; randomized, success $1/2$
& $\Omega(\mathcal P_d(m_{\rm cvx}))$
& $d\ge d_0$, $\eps\le cLR$ \\
Dorn et al.~\cite{DornEtAl2026}, Euclidean case
& Strongly convex; randomized, success $2/3$
& $\Omega(\mathcal P_d(\sigma_\eps))$
& $d\ge m_Z^\circ$ and \eqref{eq:dorn-window} \\
This paper
& Strongly convex; expected error and success $1/2$
& $\Omega(\mathcal P_d(m_\eps))$;
upper $O(d m_\eps\logfac)$
& Every $d\ge1$; $\eps\le c_0S_\Delta$;
$\Delta=0$ requires no queries \\
\bottomrule
\end{tabularx}
\end{table}

\paragraph{Organization.}
Section~\ref{sec:model} states the model and main theorem.
Section~\ref{sec:geometry} proves the slack geometry, and
Section~\ref{sec:upper} gives the algorithms.
Sections~\ref{sec:source}--\ref{sec:transfer} develop the
support-function reduction. Section~\ref{sec:spherical} handles its
finite-dimensional endpoints. Section~\ref{sec:calibration} completes
the parameter calibration and the proof of the main theorem.

\section{Oracle model and main result}\label{sec:model}

\subsection{Function class and information access}

For $R>0$, write
\[
 \B_R^d=\{x\in\R^d:\norm{x}\le R\},
 \qquad \Sph^{d-1}=\{x\in\R^d:\norm{x}=1\}.
\]
Given $d\in\mathbb N$ and $L,\mu,R>0$, let
$\cF^{\mathrm{sc}}_{L,\mu,R}$ consist of all continuous functions
$F:\B_R^d\to\R$ such that
\begin{equation}\label{eq:class}
 F-\frac{\mu}{2}\norm{\cdot}^2\ \text{is convex},
 \qquad
 |F(x)-F(y)|\le L\norm{x-y}\quad(x,y\in\B_R^d).
\end{equation}
Thus $L$ is an upper bound on the Lipschitz constant of the
\emph{whole objective}, and $\mu$ is a lower bound on its strong
convexity modulus. Neither differentiability nor an interior minimizer
is assumed. The parameters and the ball are known to the algorithm.

The oracle returns only $F(x)$ at a feasible point $x$. The unknown
objective is fixed before the algorithm's independent random seed is
drawn. Query and output maps are Borel measurable jointly in that seed
and the previous scalar responses. The seed may take values in any
standard Borel probability space. Internal computation and exact real
arithmetic are free; all scalar evaluations are charged. In particular,
a two-point estimate costs two queries. The output must lie in
$\B_R^d$ but need not have been queried.

We use a worst-case budget of at most $T$ queries on every transcript;
an algorithm can always be padded to exactly $T$ queries. Zero-query
algorithms are allowed. For the output $\widehat x_T$, set
\[
 F^*=\min_{x\in\B_R^d}F(x),
 \qquad \err_T(F)=F(\widehat x_T)-F^*.
\]
The class in \eqref{eq:class} is a closed subset of
$\mathcal C(\B_R^d)$ in the uniform norm. Evaluation and minimization
are continuous in the relevant arguments, so induction through the
query maps makes the transcript and loss jointly measurable.
Moreover, $0\le\err_T(F)\le2LR$.

\subsection{Risk criteria}

\begin{definition}[Query complexities]\label{def:risks}
For a function class $\cF$, define
\begin{align}
 T^{\mathrm{exp}}_\eps(\cF)
 &:=\inf\left\{T\in\mathbb N_0:
     \exists A\ \text{using at most }T\text{ queries},\quad
     \sup_{F\in\cF}\E_A\err_T(F)\le\eps\right\},
       \label{eq:exp-risk}\\
 T^{\mathrm{prob}}_\eps(\cF)
 &:=\inf\left\{T\in\mathbb N_0:
     \exists A\ \text{using at most }T\text{ queries},\quad
     \sup_{F\in\cF}\Pp_A(\err_T(F)>\eps)\le\frac12\right\}.
       \label{eq:prob-risk}
\end{align}
The infimum of an empty set of budgets is $+\infty$.
\end{definition}

The expectation and probability in these definitions refer only to
the algorithm's seed. Random objectives appear later solely as a
device for proving worst-case lower bounds.

\subsection{The uniform complexity theorem}

Let $\Delta=L-\mu R$. For $\Delta\ge0$, define
\begin{equation}\label{eq:S-delta}
 S_\Delta:=\max_{0\le\rho\le R}
       \left\{\Delta\rho-\frac{\mu}{2}\rho^2\right\}
 =\begin{cases}
    \Delta^2/(2\mu),&0\le\Delta\le\mu R,\\[2pt]
    \Delta R-\mu R^2/2,&\Delta>\mu R.
   \end{cases}
\end{equation}
This scale satisfies
\begin{equation}\label{eq:S-comparison}
 \frac12\min\{\Delta R,\Delta^2/\mu\}
 \le S_\Delta\le\min\{\Delta R,\Delta^2/\mu\}.
\end{equation}
For $\Delta>0$ and $\eps>0$, put
\begin{equation}\label{eq:m-eps}
 m_\eps:=\min\left\{d,\left(\frac{\Delta R}{\eps}\right)^2,
                           \frac{\Delta^2}{\mu\eps}\right\},
 \qquad
 \logfac:=\log(d+1)\log\!\left(\frac{128S_\Delta}{\eps}\right).
\end{equation}
The logarithmic factor is used only when $\eps\le S_\Delta$.
All logarithms are natural. Universal constants do not depend on
$d,L,\mu,R$, or $\eps$.

\Needspace{18\baselineskip}
\begin{theorem}[Uniform exact-value complexity]\label{thm:main}
The class $\cF^{\mathrm{sc}}_{L,\mu,R}$ is nonempty if and only if
$L\ge\mu R$. At $L=\mu R$, every member satisfies
$F(x)=F(0)+\mu\norm{x}^2/2$, and both query complexities are zero.

There exist universal $c,C,c_0>0$ such that, for every $d\ge1$,
$L>\mu R$, and $0<\eps\le c_0S_\Delta$, the following holds for
$\mathsf{risk}\in\{\mathrm{exp},\mathrm{prob}\}$:
\begin{equation}\label{eq:main-two-sided}
 c\frac{d m_\eps}{\log(e m_\eps)}
 \le T^{\mathsf{risk}}_\eps(\cF^{\mathrm{sc}}_{L,\mu,R})
 \le C\min\left\{
       \frac{d\Delta^2R^2}{\eps^2},
       \frac{d\Delta^2}{\mu\eps},
       d^2\logfac\right\}.
\end{equation}
In particular, the upper bound is at most $C d m_\eps\logfac$.
For every $\eps\ge S_\Delta$, returning the origin uses zero queries
and satisfies both criteria. The minimax expected error with zero
queries is exactly $S_\Delta$.
\end{theorem}

The proof is completed in Section~\ref{sec:calibration}.
Equation~\eqref{eq:main-two-sided} is often summarized as
\begin{equation}\label{eq:tilde-theta}
 T^{\mathsf{risk}}_\eps(\cF^{\mathrm{sc}}_{L,\mu,R})
 =\widetilde\Theta(d m_\eps),
 \qquad 0<\eps\le c_0S_\Delta.
\end{equation}
Here the suppressed factors are precisely those bounded by the
logarithms displayed in \eqref{eq:main-two-sided} and \eqref{eq:m-eps}.
The theorem determines the polynomial profile, not the optimal
logarithmic dependence. It also leaves a constant-width relative
accuracy band $c_0S_\Delta<\eps<S_\Delta$ outside the two-sided estimate.

\subsection{Accuracy regimes and consequences}

Table~\ref{tab:regimes} gives the exact inequalities selecting the
three branches of $m_\eps$. Each row is intersected with
$0<\eps\le c_0S_\Delta$, and some rows may be empty for a given
choice of parameters. The radius and strong-convexity branches meet
at $\eps=\mu R^2$; either may meet the dimension cap first.

\begin{table}[H]
\centering
\small
\setlength{\tabcolsep}{4pt}
\renewcommand{\arraystretch}{1.35}
\caption{Branches of the uniform theorem. Write
$a_\eps=(\Delta R/\eps)^2$ and $b_\eps=\Delta^2/(\mu\eps)$.
The constants in the lower and upper bounds are universal. Boundary
equalities may be assigned to either adjacent row.}
\label{tab:regimes}
\begin{tabularx}{\textwidth}{@{}P{0.17\textwidth}Y P{0.21\textwidth}P{0.15\textwidth}@{}}
\toprule
Active branch & Conditions on $\eps$ & Lower bound & Upper bound \\
\midrule
Radius, $m_\eps=a_\eps$
& $\eps\ge\max\{\mu R^2,\Delta R/\sqrt d\}$
& $\Omega(d a_\eps/\log(ea_\eps))$
& $O(d a_\eps)$ \\
Strong convexity, $m_\eps=b_\eps$
& $\Delta^2/(\mu d)\le\eps\le\mu R^2$
& $\Omega(d b_\eps/\log(eb_\eps))$
& $O(d b_\eps)$ \\
Dimension, $m_\eps=d$
& $\eps\le\min\{\Delta R/\sqrt d,\Delta^2/(\mu d)\}$
& $\Omega(d^2/\log(ed))$
& $O(d^2\logfac)$ \\
\bottomrule
\end{tabularx}
\end{table}

\begin{corollary}[Parameters separated from the boundary]\label{cor:nondegenerate}
Fix $\vartheta\in(0,1)$. If $\mu R\le(1-\vartheta)L$, then
$\vartheta L\le\Delta\le L$. Consequently, in the accuracy range
of \cref{thm:main}, the polynomial complexity profile can equivalently
be written
\[
 d\min\left\{d,\frac{L^2R^2}{\eps^2},\frac{L^2}{\mu\eps}\right\},
\]
with constants depending only on $\vartheta$.
\end{corollary}

This corollary explains when the usual $L$-based expression is valid.
It does not extend uniformly to $\vartheta\downarrow0$.

\section{Geometry forced by total Lipschitzness}\label{sec:geometry}

Set
\[
   h(x):=F(x)-\frac{\mu}{2}\norm{x}^2,
   \qquad
   \Delta:=L-\mu R.
\]
The next theorem identifies the correct uniform parameter.

\Needspace{19\baselineskip}
\begin{theorem}[Slack geometry]\label{thm:slack-geometry}
Let $F\in\cF^{\mathrm{sc}}_{L,\mu,R}$.  Then:
\begin{enumerate}[label=(\roman*),leftmargin=2.2em]
\item $L\ge\mu R$, and hence $\Delta\ge0$.
\item For every $u\in\Sph^{d-1}$, the radial residual
\[
   \psi_u(t):=h(tu),\qquad t\in[-R,R],
\]
is $\Delta$-Lipschitz.
\item The unique minimizer $x^*$ satisfies
\[
   \norm{x^*}\le \min\left\{R,\frac{\Delta}{\mu}\right\}.
\]
\item For every $0\le r<R$, the residual $h$ is
\[
   G(r):=L-\mu\sqrt{R^2-r^2}
\]
Lipschitz on $\B_r^d$.
\item If $\Delta=0$, then
\[
   F(x)=F(0)+\frac{\mu}{2}\norm{x}^2
   \qquad(x\in\B_R^d).
\]
\end{enumerate}
\end{theorem}

\begin{proof}
Fix $u\in\Sph^{d-1}$ and define $\phi_u(t)=F(tu)$ on $[-R,R]$.  The function $\phi_u$ is $L$-Lipschitz and $\mu$-strongly convex.  For $0<\eta<R$, choose
\[
  p_+\in\partial\phi_u(R-\eta),
  \qquad
  p_-\in\partial\phi_u(-R+\eta).
\]
Interior subgradients of an $L$-Lipschitz convex function lie in $[-L,L]$, while strong monotonicity gives
\[
  p_+-p_-\ge 2\mu(R-\eta).
\]
Thus $2L\ge2\mu(R-\eta)$, and $\eta\downarrow0$ proves $L\ge\mu R$.

Now $\psi_u(t)=\phi_u(t)-\mu t^2/2$ is convex.  Fix $t\in(-R,R)$ and $q\in\partial\psi_u(t)$.  For $t<r<R$, choose $q_r\in\partial\psi_u(r)$.  Monotonicity and the sum rule imply
\[
  q\le q_r,
  \qquad
  q_r+\mu r\in\partial\phi_u(r),
\]
so $q\le L-\mu r$.  Letting $r\uparrow R$ gives $q\le\Delta$.  The symmetric argument toward $-R$ gives $q\ge-\Delta$.  Hence every interior subgradient of $\psi_u$ lies in $[-\Delta,\Delta]$, proving radial $\Delta$-Lipschitzness.

Let $x^*$ be the unique minimizer and write $x^*=\rho u$ when $\rho>0$.  If $0<\rho<R$, then $0\in\partial\phi_u(\rho)$, so $-\mu\rho\in\partial\psi_u(\rho)$ and therefore $\mu\rho\le\Delta$.  If $\rho=R$, then the radial minimum occurs at the right endpoint.  For every $t<R$ and $p_t\in\partial\phi_u(t)$, the subgradient inequality gives $p_t\le0$.  Strong monotonicity between $s<t$ then yields $L\ge\mu(t-s)$.  Sending $t\uparrow R$ and $s\downarrow-R$ gives $L\ge2\mu R$, hence $\Delta\ge\mu R$ and again $R\le\Delta/\mu$.

For the local residual bound, fix $x\in\B_r^d$ and $g\in\partial h(x)$ with $g\ne0$.  Put $u=g/\norm{g}$ and let $y=x+tu$ be the first intersection of the ray $x+\R_+u$ with $\partial\B_R^d$.  Decomposing $x$ along $u$ gives
\[
   \inner{y}{u}\ge\sqrt{R^2-r^2}.
\]
For $0<\theta<1$, set $y_\theta=(1-\theta)y+\theta x$ and choose $g_\theta\in\partial h(y_\theta)$.  Monotonicity gives $\inner{g_\theta}{u}\ge\norm{g}$, while
\[
   g_\theta+\mu y_\theta\in\partial F(y_\theta)
\]
has norm at most $L$.  Therefore
\[
  L\ge \inner{g_\theta+\mu y_\theta}{u}
    \ge \norm{g}+\mu\inner{y_\theta}{u}.
\]
Letting $\theta\downarrow0$ gives $\norm{g}\le G(r)$.  Bounded subgradients on line segments imply that $h$ is $G(r)$-Lipschitz on $\B_r^d$.

Finally, if $\Delta=0$, every radial restriction $\psi_u$ is constant.  Every point in the ball lies on a radial segment through the origin, so $h(x)=h(0)$ throughout $\B_R^d$.
\end{proof}


\begin{corollary}[The exact zero-query expected-risk scale]\label{cor:zero-query}
For every $F\in\cF^{\mathrm{sc}}_{L,\mu,R}$,
$F(0)-F^*\le S_\Delta$. Moreover,
\begin{equation}\label{eq:zero-minimax}
 \inf_{A:\,T=0}\ \sup_{F\in\cF^{\mathrm{sc}}_{L,\mu,R}}
     \E_A\err_0(F)=S_\Delta.
\end{equation}
Thus $\eps\ge S_\Delta$ is the exact zero-query threshold for
the expected-error criterion.
\end{corollary}

\begin{proof}
For a minimizer $x^*=\rho u$, radial Lipschitzness gives
\[
 F(0)-F(x^*)=h(0)-h(\rho u)-\frac{\mu}{2}\rho^2
 \le\Delta\rho-\frac{\mu}{2}\rho^2\le S_\Delta.
\]
The case $x^*=0$ is immediate. For the reverse minimax inequality,
let $V$ be uniform on $\Sph^{d-1}$ and consider
\[
 Q_V(x)=\frac{\mu}{2}\norm{x}^2-\Delta\inner{V}{x}.
\]
Every realization is admissible and has minimum $-S_\Delta$.
Any zero-query output $X$ is independent of $V$, so symmetry yields
\[
 \E_{V,A}[Q_V(X)+S_\Delta]
 =S_\Delta+\frac{\mu}{2}\E_A\norm{X}^2\ge S_\Delta.
\]
At least one fixed realization attains this average lower bound.
Returning zero attains the upper bound.
\end{proof}

The exact threshold in \eqref{eq:zero-minimax} concerns expected
error. We use it as a normalization for both criteria, without
asserting the same exact zero-query threshold for success probability
$1/2$.

\section{Localized upper bounds}\label{sec:upper}

The upper bound uses standard value-oracle methods after a geometric
localization. We give the symmetric estimator of Shamir~\cite{Shamir2017}
and the projected-iteration estimates
explicitly to separate their guarantees from the localization step.

\subsection{Two-point estimation and projected iterations}

\begin{lemma}[Symmetric two-point estimator]\label{lem:two-point}
Let $h$ be convex and $G$-Lipschitz on $\B_S^d$. For
$0<2\delta<S$, set
\[
 X_\delta=\left(1-\frac{2\delta}{S}\right)\B_S^d,
 \qquad
 h_\delta(x)=\E_{W\sim\mathrm{Unif}(\B_1^d)}h(x+\delta W).
\]
For $U\sim\mathrm{Unif}(\Sph^{d-1})$, define
\[
 \widehat g_{h,\delta}(x,U)=
 \frac{d}{2\delta}\bigl(h(x+\delta U)-h(x-\delta U)\bigr)U.
\]
Then $h_\delta$ is convex and continuously differentiable on a
neighborhood of $X_\delta$, and
\begin{equation}\label{eq:two-point-properties}
 \E_U\widehat g_{h,\delta}(x,U)=\nabla h_\delta(x),\qquad
 |h_\delta(x)-h(x)|\le G\delta,\qquad
 \E_U\norm{\widehat g_{h,\delta}(x,U)}^2\le2dG^2.
\end{equation}
\end{lemma}

\begin{proof}
The ball average preserves convexity, and Lipschitzness gives the
bias bound. Differentiation of the ball integral and symmetry give
\[
 \nabla h_\delta(x)=\frac d\delta\E_U[h(x+\delta U)U]
                  =\E_U\widehat g_{h,\delta}(x,U).
\]
For the second moment, the spherical Poincar\'e inequality gives, for $d\ge2$ and
$q(u)=h(x+\delta u)$,
\[
 \operatorname{Var}(q(U))
 \le\frac{1}{d-1}\E\|\nabla_{\Sph}q(U)\|_2^2
 \le\frac{G^2\delta^2}{d-1}.
\]
The inequality holds for Lipschitz $q$ by approximation. Since
$q(U)$ and $q(-U)$ have the same mean,
\[
 \E(q(U)-q(-U))^2\le4\operatorname{Var}(q(U)).
\]
Multiplication by $d^2/(4\delta^2)$ proves the claimed bound because
$d^2/(d-1)\le2d$. For $d=1$, Lipschitzness directly bounds the
estimator by $G$. All points used in the argument lie inside $\B_S^d$.
\end{proof}

\begin{lemma}[Projected iteration bounds]\label{lem:projected}
Let $X$ be a closed convex subset of $\B_S^d$ containing zero,
and let $\phi$ be continuous and convex on $X$. Suppose $g_t$ is conditionally
unbiased for a subgradient at $x_t$ and
$\E_t\norm{g_t}^2\le\Gamma^2$. Starting at $x_1=0$, use
\[
 x_{t+1}=\operatorname{Proj}_X(x_t-\eta_tg_t).
\]
For constant $\eta_t=S/(\Gamma\sqrt N)$, the average
$\bar x_N=N^{-1}\sum_{t=1}^Nx_t$ satisfies
\[
 \E\phi(\bar x_N)-\min_X\phi\le S\Gamma/\sqrt N.
\]
If $\phi$ is $\mu$-strongly convex, then
$\eta_t=2/[\mu(t+1)]$ and
$\bar x_N=2\sum_{t=1}^Ntx_t/[N(N+1)]$ give
\[
 \E\phi(\bar x_N)-\min_X\phi\le\frac{2\Gamma^2}{\mu(N+1)}.
\]
\end{lemma}

\begin{proof}
Let $z$ minimize $\phi$ on $X$ and write
$D_t=\E\norm{x_t-z}^2$. Nonexpansiveness of projection and the
subgradient inequality give
\[
 \E[\phi(x_t)-\phi(z)]
 \le\frac{D_t-D_{t+1}}{2\eta_t}+\frac{\eta_t\Gamma^2}{2}.
\]
Summation, $D_1\le S^2$, and convexity prove the first estimate.
Under strong convexity, the right side has the additional term
$-\mu D_t/2$. With the stated decreasing step size, multiplication
by $t$ yields
\[
 t\E[\phi(x_t)-\phi(z)]
 \le\frac{\mu t(t-1)}4D_t
     -\frac{\mu t(t+1)}4D_{t+1}
     +\frac{t\Gamma^2}{\mu(t+1)}.
\]
The distance terms telescope, and the remaining sum is at most
$N\Gamma^2/\mu$. Division by $N(N+1)/2$ and convexity finish the proof.
\end{proof}

\subsection{Localization and the three algorithmic branches}

\begin{theorem}[Localized exact-value upper bound]\label{thm:upper}
There is a universal $C>0$ such that, for every $\Delta>0$ and
$0<\eps\le S_\Delta$, both risk criteria satisfy
\begin{equation}\label{eq:upper-envelope}
 T^{\mathsf{risk}}_\eps(\cF^{\mathrm{sc}}_{L,\mu,R})
 \le C\min\left\{
       \frac{d\Delta^2R^2}{\eps^2},
       \frac{d\Delta^2}{\mu\eps},
       d^2\log(d+1)\log\!\left(\frac{128S_\Delta}{\eps}\right)
       \right\}.
\end{equation}
\end{theorem}

\begin{proof}
Write $F=h+\mu\norm{\cdot}^2/2$ and choose the known working radius
\begin{equation}\label{eq:working-radius}
 S=\begin{cases}
      2\Delta/\mu,&0<\Delta<\mu R/2,\\
      R,&\Delta\ge\mu R/2.
    \end{cases}
\end{equation}
By \cref{thm:slack-geometry}, a minimizer lies in $\B_S^d$.
In the first case, $S<R$ and the local estimate gives
\[
 \Lip_{\B_S^d}(h)
 \le\Delta+\frac{\mu S^2}{R+\sqrt{R^2-S^2}}
 \le\Delta+\frac{4\Delta^2}{\mu R}<3\Delta.
\]
In the second case,
$\Lip_{\B_R^d}(h)\le L+\mu R=\Delta+2\mu R\le5\Delta$.
Thus, in both cases,
\begin{equation}\label{eq:localized-G}
 \Lip_{\B_S^d}(h)\le\bar G:=5\Delta,
 \qquad \mu S\le2\Delta,
 \qquad S_\Delta\le\Delta S\le4S_\Delta.
\end{equation}
The last inequalities follow directly from \eqref{eq:S-delta} and
\eqref{eq:working-radius}.

Set $\delta=\eps/(32\bar G)$ and use the set $X_\delta$ of
\cref{lem:two-point}. Equation~\eqref{eq:localized-G} ensures
$2\delta<S$. Optimize
\[
 \Phi_\delta(x)=h_\delta(x)+\frac{\mu}{2}\norm{x}^2
\]
on $X_\delta$. Draw $U_t$ independently and uniformly on the unit
sphere at each iteration. Two exact values of $F$ at $x_t\pm\delta U_t$ reveal
the corresponding residual values after subtracting the known
quadratic. The estimator
\[
 g_t=\widehat g_{h,\delta}(x_t,U_t)+\mu x_t
\]
is unbiased for $\nabla\Phi_\delta(x_t)$ and satisfies
\begin{equation}\label{eq:gradient-moment}
 \E_t\norm{g_t}^2
 \le4d\bar G^2+2\mu^2S^2\le108d\Delta^2=:\Gamma^2.
\end{equation}

Let $x^*$ minimize $F$ and put $z=(1-2\delta/S)x^*$.
Then $z\in X_\delta$, $\norm{z-x^*}\le2\delta$, and its quadratic
term is no larger than that of $x^*$. Therefore
$\Phi_\delta(z)-F(x^*)\le3\bar G\delta$.
For any output $\bar x\in X_\delta$,
\begin{equation}\label{eq:smoothing-total}
 F(\bar x)-F^*
 \le\Phi_\delta(\bar x)-\min_{X_\delta}\Phi_\delta
       +4\bar G\delta.
\end{equation}
Thus the smoothing and contraction cost is at most $\eps/8$.

The two schedules in \cref{lem:projected}, run to optimization
error at most $\eps/4$, require respectively
\[
 N=O\left(\frac{d\Delta^2S^2}{\eps^2}\right)
   \le O\left(\frac{d\Delta^2R^2}{\eps^2}\right),
 \qquad
 N=O\left(\frac{d\Delta^2}{\mu\eps}\right).
\]
Each iteration uses two feasible evaluations. The resulting expected
error is below $\eps/2$, so Markov's inequality also gives success
probability at least $1/2$ at accuracy $\eps$.

For the final branch, $F$ is $M$-Lipschitz on $\B_S^d$ with
$M=7\Delta$, by \eqref{eq:localized-G}. Protasov's deterministic
value-only bound~\cite{Protasov1996}, applied to this known ball,
requires
\[
 O\left(d^2\log(d+1)\log\!\frac{4MS}{\eps}\right)
 \le O\left(d^2\log(d+1)\log\!\frac{128S_\Delta}{\eps}\right)
\]
evaluations. Here $4MS\le112S_\Delta$ and
$\eps\le S_\Delta\le MS$. The minimizer over $\B_S^d$ is also
a minimizer over the original ball. This deterministic guarantee
implies both risk criteria. Taking the best of the three algorithms
proves \eqref{eq:upper-envelope}.
\end{proof}

The new ingredient in this upper bound is the localization under the
total Lipschitz constraint. The directional estimator, projected
updates, and general value-only algorithm are established methods.

\section{The support-function source}\label{sec:source}

Let $B_1,\ldots,B_k$ be independent standard Gaussian vectors conditioned on
\[
  \norm{B_i}\le2\sqrt d,
\]
and set
\[
  a_i:=\frac{B_i}{\sqrt d},
  \qquad
  h_\Xi(z):=\max_{1\le i\le k}\inner{a_i}{z},
  \qquad z\in\R^d.
\]
The formula defines the function on all of $\R^d$; the source oracle
accepts queries only in $\B_1^d$. Every realization is convex,
positively homogeneous, and $2$-Lipschitz. The induced law on
$\mathcal C(\B_1^d)$ is a fixed Borel probability measure depending
only on $(d,k)$.

The following two results are the external information-theoretic
input. They concern the same prior, which is essential for the
intersection of events used below. We cite their precise fixed-prior
form, rather than only the resulting minimax bound.

\begin{theorem}[Source facts: {\cite[Theorems~5 and~40]{ZhangEtAl2026}}]\label{thm:source-facts}
There are universal constants $c_{\mathrm{dim}},c_{\mathrm{qry}},c_{\mathrm{src}}>0$
and $d_{\mathrm{src}}\in\mathbb N$ such that, with
$c_{\mathrm{sv}}=1/4$, whenever
\[
   d\ge d_{\mathrm{src}},
   \qquad
   1\le k\le c_{\mathrm{dim}}d,
\]
the following hold for the common source law above.

\begin{enumerate}[label=(\roman*),leftmargin=2.2em]
\item For every deterministic adaptive exact-value source algorithm using
\[
   T\le c_{\mathrm{qry}}\frac{dk}{\log(ek)}
\]
queries,
\begin{equation}\label{eq:source-lb}
  \Pp\left(
      h_\Xi(\widehat z_T)-\min_{z\in\B_1^d}h_\Xi(z)
      >\frac{c_{\mathrm{src}}}{\sqrt k}
      \right)\ge\frac34.
\end{equation}
The same source prior is used for every deterministic algorithm.

\item Let $A_\Xi\in\R^{k\times d}$ have rows $a_i^\top$.  Then
\begin{equation}\label{eq:spectral-event}
  \Pp\left(A_\Xi A_\Xi^\top\succeq c_{\mathrm{sv}}I_k\right)\ge\frac78.
\end{equation}
\end{enumerate}
\end{theorem}

The common-prior clause in \eqref{eq:source-lb} is essential: it permits different source algorithms to be tested on the same seed space and permits integration over an independent random tape.


Theorem~40 of the cited source applies when $k\le d/16$, so we
decrease $c_{\mathrm{dim}}$ if necessary. The zero-query case in
part~(i) follows by padding with one unused query after increasing
$d_{\mathrm{src}}$ so that $c_{\mathrm{qry}}d_{\mathrm{src}}\ge1$.
No posterior-information estimate from that source is reproved or
claimed as a contribution here.

\section{Quadratic regularization and radial optimization}\label{sec:radial}

Fix $\Delta>0$ and put
\[
   \alpha:=\frac{\Delta}{2}.
\]
For a source realization $\xi$, define the strongly convex target
\begin{equation}\label{eq:hard-family}
   F_\xi(x):=\alpha h_\xi(x)+\frac{\mu}{2}\norm{x}^2,
   \qquad x\in\B_R^d.
\end{equation}
Because $h_\xi$ is $2$-Lipschitz, $F_\xi$ is $L=(\Delta+\mu R)$-Lipschitz and $\mu$-strongly convex.  The seed-to-function map is continuous in the uniform norm, so the induced target law is Borel.

\begin{proposition}[Exact preservation of scalar queries]\label{prop:conjugacy}
One scalar query to $F_\xi$ is simulated by one scalar query to $h_\xi$ and conversely.  At $x=Rz$,
\begin{equation}\label{eq:conjugacy}
  F_\xi(Rz)=\alpha R h_\xi(z)+\frac{\mu R^2}{2}\norm{z}^2,
  \qquad
  h_\xi(z)=\frac{F_\xi(Rz)-\mu R^2\norm{z}^2/2}{\alpha R}.
\end{equation}
The transformations preserve adaptive Borel transcripts and charge exactly one source query per target query.
\end{proposition}

\begin{proof}
The identities use positive homogeneity of $h_\xi$ and $\alpha R>0$.  Given a target history, subtract the known quadratic and divide by the known scale; given a source history, reverse the arithmetic.  Induction through the deterministic query maps gives mutually inverse Borel transformations of the complete scalar histories.
\end{proof}

For $q\ge0$, define the profile
\begin{equation}\label{eq:Psi-def}
  \Psi(q):=\max_{0\le\rho\le R}
     \left\{\alpha q\rho-\frac{\mu}{2}\rho^2\right\}.
\end{equation}
It has the explicit form
\begin{equation}\label{eq:Psi-form}
  \Psi(q)=
  \begin{cases}
    \dfrac{\alpha^2q^2}{2\mu},&\alpha q\le\mu R,\\[5pt]
    \alpha Rq-\dfrac{\mu R^2}{2},&\alpha q>\mu R.
  \end{cases}
\end{equation}
The map is convex, increasing, and vanishes at zero.  Moreover,
\begin{equation}\label{eq:Psi-lower}
  \Psi(q)\ge\frac12\min\left\{\alpha Rq,\frac{\alpha^2q^2}{\mu}\right\},
\end{equation}
and, for $0\le a\le q$,
\begin{equation}\label{eq:Psi-increment}
  \Psi(q)-\Psi(q-a)\ge\Psi(a).
\end{equation}
Indeed, increments of a convex function over intervals of fixed length are nondecreasing.

Define the negative source depth
\[
   q_\xi:=-\min_{z\in\B_1^d}h_\xi(z)\ge0.
\]

\begin{lemma}[Radial variational formula]\label{lem:radial}
For every source realization,
\begin{equation}\label{eq:exact-optimum}
   \min_{x\in\B_R^d}F_\xi(x)=-\Psi(q_\xi).
\end{equation}
If $x=\rho u$ with $0\le\rho\le R$, $\norm{u}=1$, and
\[
   y:=(-h_\xi(u))_+,
\]
then
\begin{equation}\label{eq:radial-gap}
   F_\xi(x)-\min F_\xi\ge \Psi(q_\xi)-\Psi(y).
\end{equation}
\end{lemma}

\begin{proof}
If $q_\xi=0$, then $h_\xi\ge0$ on the unit ball and the optimum is zero at the origin.  Suppose $q_\xi>0$.  Positive homogeneity implies that any minimizer of $h_\xi$ lies on the unit sphere: otherwise radial expansion would make its negative value more negative.  Hence there is $u^*$ with $h_\xi(u^*)=-q_\xi$.  Along the ray $\rho u^*$, minimizing $F_\xi$ is exactly the one-dimensional problem defining $-\Psi(q_\xi)$, and no other direction can do better because $h_\xi(u)\ge-q_\xi$.

For the reported direction, $h_\xi(u)\ge-y$, so
\[
  F_\xi(\rho u)\ge-\alpha y\rho+\frac{\mu}{2}\rho^2\ge-\Psi(y).
\]
Subtract \eqref{eq:exact-optimum}.
\end{proof}

The depth has an equivalent matrix representation.

\begin{lemma}[Spectral depth]\label{lem:spectral-depth}
Let $A_\xi$ have rows $a_i^\top$.  Then
\begin{equation}\label{eq:q-simplex}
   q_\xi=\min_{p\in\Delta_k}\norm{A_\xi^\top p},
   \qquad
   \Delta_k:=\{p\in\R_+^k:\mathbf 1^\top p=1\}.
\end{equation}
Consequently, on the event $A_\Xi A_\Xi^\top\succeq c_{\mathrm{sv}}I_k$,
\begin{equation}\label{eq:q-spectral}
   q_\Xi\ge\sqrt{\frac{c_{\mathrm{sv}}}{k}}.
\end{equation}
\end{lemma}

\begin{proof}
Let $K=\operatorname{conv}\{a_1,\ldots,a_k\}$ and let $b$ be its
point of minimum norm. If $b=0$, then $h_\xi(z)\ge0$ for every $z$,
so both sides of \eqref{eq:q-simplex} vanish. Otherwise, the
projection inequality gives $\inner{a_i-b}{b}\ge0$ for every $i$.
Consequently,
\[
 h_\xi(-b/\norm{b})\le-\norm{b}.
\]
Conversely, $h_\xi(z)\ge\inner{b}{z}\ge-\norm{b}$ on the unit
ball. Thus $q_\xi=\norm{b}=\min_{p\in\Delta_k}\norm{A_\xi^\top p}$.
On the spectral event,
\[
 \norm{A_\xi^\top p}^2
 =p^\top A_\xi A_\xi^\top p
 \ge c_{\mathrm{sv}}\norm{p}^2
 \ge c_{\mathrm{sv}}/k,
\]
which proves \eqref{eq:q-spectral}.
\end{proof}

\section{Transferring source hardness}\label{sec:transfer}

Fix a deterministic $T$-query target algorithm for the family \eqref{eq:hard-family}.  Use \cref{prop:conjugacy} to simulate its queries on the source.  Let $X_\xi$ be its target output and define the Borel direction map
\[
   U(x):=
   \begin{cases}
      x/\norm{x},&x\ne0,\\
      u_0,&x=0,
   \end{cases}
\]
for a fixed $u_0\in\Sph^{d-1}$.

Consider two source algorithms with the same simulated queries and scalar transcript:
\begin{itemize}[leftmargin=2em]
\item the \emph{zero-output algorithm} always outputs $0$;
\item the \emph{direction-output algorithm} outputs $U(X_\xi)$.
\end{itemize}
The zero-output algorithm has source error $q_\xi$, while the direction-output algorithm has source error
\[
   h_\xi(U(X_\xi))+q_\xi.
\]

\begin{proposition}[Expected-error transfer]\label{prop:expected-transfer}
Under the assumptions and query budget of \cref{thm:source-facts}, every deterministic target algorithm satisfies
\begin{equation}\label{eq:expected-fixed-budget}
  \E_\Xi\bigl[F_\Xi(X_\Xi)-\min F_\Xi\bigr]
  \ge \frac12\Psi\!\left(\frac{c_{\mathrm{src}}}{\sqrt k}\right)
  \ge c\min\left\{\frac{\Delta R}{\sqrt k},
                    \frac{\Delta^2}{\mu k}\right\}.
\end{equation}
\end{proposition}

\begin{proof}
Apply \eqref{eq:source-lb} separately to the zero-output and direction-output algorithms.  On the same source seed space, the events
\[
  E_0:=\left\{q_\Xi>\frac{c_{\mathrm{src}}}{\sqrt k}\right\},
  \qquad
  E_{\mathrm{dir}}:=\left\{
       h_\Xi(U(X_\Xi))+q_\Xi>\frac{c_{\mathrm{src}}}{\sqrt k}
       \right\}
\]
each have probability at least $3/4$, hence their intersection has probability at least $1/2$.  Put $\eta=c_{\mathrm{src}}/\sqrt k$ and $y=(-h_\Xi(U(X_\Xi)))_+$.  On the intersection, $q_\Xi>\eta$ and $y<q_\Xi-\eta$.  Therefore \cref{lem:radial} and \eqref{eq:Psi-increment} give
\[
  F_\Xi(X_\Xi)-\min F_\Xi
  >\Psi(q_\Xi)-\Psi(q_\Xi-\eta)
  \ge\Psi(\eta).
\]
The gap is nonnegative elsewhere.  Integrate and use \eqref{eq:Psi-lower} with $\alpha=\Delta/2$.
\end{proof}

The preceding argument uses only the source optimization theorem.
Its two failure events have an intersection of probability at least
$1/2$, which is enough for expected error. For the success criterion
in \eqref{eq:prob-risk}, we need a failure probability strictly larger
than $1/2$. The spectral event provides the required margin.

\begin{proposition}[Constant-success transfer]\label{prop:prob-transfer}
Under the assumptions and query budget of \cref{thm:source-facts}, every deterministic target algorithm satisfies
\begin{equation}\label{eq:prob-fixed-budget}
  \Pp_\Xi\left(
    F_\Xi(X_\Xi)-\min F_\Xi
    \ge c\min\left\{\frac{\Delta R}{\sqrt k},
                      \frac{\Delta^2}{\mu k}\right\}
  \right)\ge\frac58.
\end{equation}
\end{proposition}

\begin{proof}
Let
\[
  \gamma:=\sqrt{c_{\mathrm{sv}}},
  \qquad
  c_*:=\min\left\{c_{\mathrm{src}},\frac{\gamma}{2}\right\},
  \qquad
  \eta:=\frac{c_*}{\sqrt k}.
\]
By \cref{thm:source-facts,lem:spectral-depth}, the event
\[
  E_{\mathrm{depth}}:=\left\{q_\Xi\ge\frac{\gamma}{\sqrt k}\right\}
\]
has probability at least $7/8$.  By applying the source theorem to the direction-output algorithm,
\[
  E_{\mathrm{dir}}:=\left\{
     h_\Xi(U(X_\Xi))+q_\Xi>\frac{c_{\mathrm{src}}}{\sqrt k}
  \right\}
\]
has probability at least $3/4$.  Hence
\[
  \Pp(E_{\mathrm{depth}}\cap E_{\mathrm{dir}})
  \ge\frac78+\frac34-1=\frac58.
\]
On this intersection, $q_\Xi\ge2\eta$ and, with $y=(-h_\Xi(U(X_\Xi)))_+$, one has $y<q_\Xi-\eta$.  The same radial increment argument gives a target gap at least $\Psi(\eta)$.  Apply \eqref{eq:Psi-lower}.
\end{proof}

\begin{remark}[Information and geometry have separate roles]
The source theorem controls the quality of the reported direction;
the spectral event controls the depth of the best direction. Their
intersection is estimated by a union bound, with no independence
assumption. Quadratic regularization changes the conversion from
directional error to objective error through $\Psi$, while preserving
the scalar transcript exactly.
\end{remark}

\section{Finite-dimensional endpoints}\label{sec:spherical}

When $m_\eps$ is bounded, the target lower bound is only of order
$d$. A simpler family supplies this endpoint and avoids any lower
threshold on the ambient dimension.

For $v\in\Sph^{d-1}$, define
\begin{equation}\label{eq:spherical-family}
   Q_v(x):=\frac{\mu}{2}\norm{x}^2-\Delta\inner{v}{x},
   \qquad x\in\B_R^d.
\end{equation}
Every $Q_v$ belongs to $\cF^{\mathrm{sc}}_{L,\mu,R}$, since
\[
  \norm{\nabla Q_v(x)}\le\mu R+\Delta=L.
\]
Its minimizer is
\begin{equation}\label{eq:spherical-minimizer}
  x_v^*=\rho_*v,
  \qquad
  \rho_*:=\min\left\{R,\frac{\Delta}{\mu}\right\}.
\end{equation}
One exact value is equivalent to one adaptive linear measurement of $v$:
\begin{equation}\label{eq:linear-measurement}
  \inner{v}{x}=\frac{\mu\norm{x}^2/2-Q_v(x)}{\Delta}.
\end{equation}
Moreover,
\[
   Q_v(0)-Q_v(x_v^*)=S_\Delta.
\]

\subsection{Posterior energy and expected risk}

\begin{lemma}[Adaptive spherical posterior energy]\label{lem:spherical-energy}
Let $V$ be uniform on $\Sph^{d-1}$ and independent of an algorithmic random tape.  Suppose an adaptive strategy observes $T\le d$ exact inner products with $V$.  If $\mathcal H_T$ is generated by the tape, queries, and responses, and
\[
   M_T:=\E[V\mid\mathcal H_T],
\]
then
\begin{equation}\label{eq:spherical-energy}
   \E\norm{M_T}^2\le\frac{T}{d}.
\end{equation}
\end{lemma}

\begin{proof}
We enlarge the observation history so that exactly one new orthogonal
coordinate is revealed at each step.  This avoids any assumption that the
rank of the original adaptive queries is deterministic.

Condition on the independent random tape, and simulate the original
strategy.  Suppose that the orthonormal vectors $E_1,\ldots,E_{t-1}$ and
coordinates $C_j=\langle V,E_j\rangle$ have already been observed.  The
original query $x_t$ is determined by its previously simulated responses.
Let
\[
  w_t=x_t-\sum_{j<t}\langle x_t,E_j\rangle E_j.
\]
If $w_t\ne0$, choose $E_t=w_t/\|w_t\|_2$.  Otherwise, choose any Borel
unit-vector completion of $E_1,\ldots,E_{t-1}$; for example, orthogonalize
the first standard basis vector with a nonzero residual.  Observe
$C_t=\langle V,E_t\rangle$.  The original response is then recovered as
\[
  \langle V,x_t\rangle
  =\sum_{j<t}\langle x_t,E_j\rangle C_j+\|w_t\|_2 C_t.
\]
Thus the enlarged experiment uses one linear measurement per original
query and reproduces the entire original transcript.  The extra
coordinate observed when $w_t=0$ is ignored by the simulated strategy.

Let $\widetilde{\mathcal H}_t$ be the enlarged history, including the
tape, and put
\[
  \widetilde M_t:=\sum_{j=1}^t C_jE_j.
\]
Conditional rotational invariance gives
$\widetilde M_t=\mathbb E[V\mid\widetilde{\mathcal H}_t]$.
More explicitly, given $\widetilde{\mathcal H}_{t-1}$, the unrevealed
component is uniform on the sphere of squared radius
$1-\|\widetilde M_{t-1}\|_2^2$ in a $(d-t+1)$-dimensional orthogonal
complement.  Hence, for $1\le t\le d$,
\[
  \mathbb E[C_t^2\mid\widetilde{\mathcal H}_{t-1}]
  =\frac{1-\|\widetilde M_{t-1}\|_2^2}{d-t+1}.
\]
Writing $a_t=\mathbb E\|\widetilde M_t\|_2^2$ yields
\[
  a_0=0,\qquad
  a_t=a_{t-1}+\frac{1-a_{t-1}}{d-t+1}=\frac{t}{d},
\]
where the final equality follows by induction.  Finally,
$\mathcal H_T\subseteq\widetilde{\mathcal H}_T$, so the tower property
and conditional Jensen inequality imply
\[
  \mathbb E\|M_T\|_2^2
  =\mathbb E\bigl\|
      \mathbb E[\widetilde M_T\mid\mathcal H_T]
    \bigr\|_2^2
  \le\mathbb E\|\widetilde M_T\|_2^2
  =\frac{T}{d}.
\]
\end{proof}

\begin{proposition}[Unified spherical Bayes-risk curve]\label{prop:spherical-expected}
For every randomized $T$-query exact-value algorithm with $0\le T\le d$, there is a fixed $v\in\Sph^{d-1}$ such that
\begin{equation}\label{eq:spherical-Bayes}
   \E\bigl[Q_v(\widehat X_T)-Q_v(x_v^*)\bigr]
   \ge S_\Delta\left(1-\sqrt{\frac{T}{d}}\right).
\end{equation}
In particular, every $T\le d/4$ algorithm has a fixed objective with expected gap at least $S_\Delta/2$.
\end{proposition}

\begin{proof}
Let $V$ be uniform and run the algorithm on $Q_V$.  By \eqref{eq:linear-measurement}, its value transcript is equivalent to an adaptive inner-product transcript, so \cref{lem:spherical-energy} applies.  Let $X$ be the output and $M=M_T$.

If $\Delta\le\mu R$, then $x_V^*=(\Delta/\mu)V$ and
\[
  Q_V(X)-Q_V(x_V^*)
  =\frac{\mu}{2}\left\|X-\frac{\Delta}{\mu}V\right\|^2.
\]
Conditioning on $\mathcal H_T$ and using the bias--variance identity gives
\[
  \E[Q_V(X)-Q_V(x_V^*)\mid\mathcal H_T]
  \ge S_\Delta(1-\norm{M}^2).
\]
Thus the average risk is at least $S_\Delta(1-T/d)$, which is no smaller than the right side of \eqref{eq:spherical-Bayes}.

If $\Delta\ge\mu R$, then $x_V^*=RV$ and the exact decomposition
\begin{equation}\label{eq:boundary-gap-decomp}
 Q_V(X)-Q_V(RV)
 =\frac{\mu}{2}\norm{X-RV}^2
   +(\Delta-\mu R)\bigl(R-\inner{V}{X}\bigr)
\end{equation}
holds.  Conditionally,
\[
  \E\norm{X-RV}^2
  =\norm{X-RM}^2+R^2(1-\norm{M}^2)
  \ge R^2(1-\norm{M}^2),
\]
and
\[
  R-\inner{M}{X}\ge R(1-\norm{M}).
\]
Using \cref{lem:spherical-energy} and Cauchy--Schwarz yields
\[
  \E[Q_V(X)-Q_V(RV)]
  \ge\left(\frac{\mu R^2}{2}+(\Delta-\mu R)R\right)
       \left(1-\sqrt{\frac{T}{d}}\right)
  =S_\Delta\left(1-\sqrt{\frac{T}{d}}\right).
\]
The left side is the average over fixed $v$, so at least one $v$ attains the same lower bound.
\end{proof}

\subsection{Constant-probability direction estimation}

\begin{lemma}[Direction estimation from few adaptive measurements]\label{lem:direction-estimation}
Let $d\ge2$, let $V$ be uniform on $\Sph^{d-1}$, and let a deterministic adaptive strategy observe $T$ exact inner products with $V$ before outputting a unit vector $U$.  If
\[
   T\le\frac{d}{256},
\]
then
\begin{equation}\label{eq:direction-prob}
   \Pp\bigl(\inner{V}{U}\le1/2\bigr)\ge\frac23.
\end{equation}
\end{lemma}

\begin{proof}
Use the enlarged experiment in the proof of
\cref{lem:spherical-energy}, with history
$\widetilde{\mathcal H}_T$ and posterior mean
$M=\E[V\mid\widetilde{\mathcal H}_T]$. The original output $U$
is measurable with respect to this history. Exactly $T$ orthogonal
coordinates have been revealed, so the remaining component is
rotationally invariant in dimension $d-T$. Consequently,
\begin{equation}\label{eq:residual-variance}
  \E\bigl[\inner{V-M}{U}^2\mid\widetilde{\mathcal H}_T\bigr]
  \le\frac{1}{d-T}.
\end{equation}
The enlarged experiment has $\E\norm{M}^2=T/d$, so Markov's inequality gives
\[
  \Pp(\norm{M}>1/4)\le16\frac{T}{d}\le\frac1{16}.
\]
If $d\ge256$, then by \eqref{eq:residual-variance},
\[
  \Pp\bigl(\inner{V-M}{U}>1/4\bigr)
  \le\frac{16}{d-T}\le\frac{16}{255}.
\]
Hence
\[
  \Pp(\inner{V}{U}>1/2)
  \le\frac1{16}+\frac{16}{255}<\frac13.
\]

If $2\le d<256$, the budget condition forces $T=0$, so $U$ is fixed.  Writing $V=G/\norm{G}$ for a standard Gaussian $G$, the event $\inner{V}{U}>1/2$ is nonincreasing in $d$ after rotating $U$ to the first coordinate: it is
\[
  \{G_1>0,\;3G_1^2>G_2^2+\cdots+G_d^2\}.
\]
At $d=2$, its angular probability is exactly $1/3$, and adding Gaussian coordinates can only decrease it.  This proves \eqref{eq:direction-prob}.
\end{proof}

\begin{lemma}[Bad direction forces objective error]\label{lem:spherical-gap}
Let $x\in\B_R^d$.  If $x=0$, or if $x\ne0$ and
\[
  \inner{v}{x/\norm{x}}\le\frac12,
\]
then
\begin{equation}\label{eq:spherical-gap}
  Q_v(x)-Q_v(x_v^*)\ge\frac38 S_\Delta.
\end{equation}
\end{lemma}

\begin{proof}
At $x=0$, the gap is exactly $S_\Delta$.  Suppose $x\ne0$ and put $u=x/\norm{x}$, $c=\inner{v}{u}\le1/2$.

If $\Delta\le\mu R$, let $\rho_*=\Delta/\mu$.  Since
\[
   Q_v(x)-Q_v(\rho_*v)=\frac{\mu}{2}\norm{x-\rho_*v}^2,
\]
the smallest distance from $\rho_*v$ to a ray whose direction has correlation at most $1/2$ is at least $\sqrt{3}\rho_*/2$.  Thus the gap is at least
\[
   \frac{3\mu\rho_*^2}{8}=\frac34S_\Delta.
\]

If $\Delta\ge\mu R$, use \eqref{eq:boundary-gap-decomp}.  Minimizing the first term over the radius on a direction with $c\le1/2$ gives
\[
   \norm{x-Rv}^2\ge\frac34R^2,
\]
while $\inner{v}{x}\le R/2$ gives $R-\inner{v}{x}\ge R/2$.  Hence
\[
  Q_v(x)-Q_v(Rv)
  \ge \frac{3\mu R^2}{8}+\frac{(\Delta-\mu R)R}{2}
  \ge\frac38\left(\Delta R-\frac{\mu R^2}{2}\right)
  =\frac38S_\Delta.
\]
\end{proof}

\begin{proposition}[Spherical constant-success endpoint, $d\ge2$]\label{prop:spherical-prob}
Let $d\ge2$.  Every deterministic algorithm using $T\le d/256$ exact-value queries satisfies, under $V\sim\mathrm{Unif}(\Sph^{d-1})$,
\begin{equation}\label{eq:spherical-prob-gap}
  \Pp\left(Q_V(\widehat X_T)-Q_V(x_V^*)\ge\frac38S_\Delta\right)
  \ge\frac23.
\end{equation}
\end{proposition}

\begin{proof}
The exact values reveal only adaptive inner products with $V$.  Apply \cref{lem:direction-estimation} to the normalized output, using an arbitrary fixed direction when the output is zero, and then apply \cref{lem:spherical-gap}.
\end{proof}

The one-dimensional constant-success criterion has a minor endpoint peculiarity because the sphere has only two points.  A three-function packing resolves it.

\begin{lemma}[One-dimensional zero-query packing]\label{lem:one-dimensional}
Let $d=1$ and $\Delta>0$.  Define on $[-R,R]$
\[
  Q_+(x)=\frac{\mu}{2}x^2-\Delta x,
  \qquad
  Q_-(x)=\frac{\mu}{2}x^2+\Delta x,
  \qquad
  Q_0(x)=\frac{\mu}{2}x^2+\Delta|x|.
\]
All three functions lie in $\cF^{\mathrm{sc}}_{L,\mu,R}$.  If $0<\eps\le S_\Delta/64$, their $\eps$-optimal sets are pairwise disjoint.  Consequently every randomized zero-query algorithm has one of these three objectives on which
\[
  \Pp(\err_0>\eps)\ge\frac23.
\]
\end{lemma}

\begin{proof}
The residuals $-\Delta x$, $\Delta x$, and $\Delta|x|$ are convex and $\Delta$-Lipschitz, so total Lipschitzness is $\Delta+\mu R=L$.

If $\Delta\le\mu R$, put $r_*=\Delta/\mu$.  The minimizers of $Q_+,Q_-,Q_0$ are $r_*,-r_*,0$.  Since $S_\Delta=\mu r_*^2/2$, $\eps\le S_\Delta/64$ implies that the $Q_+$-good set lies inside $[7r_*/8,R]$, the $Q_-$-good set lies inside $[-R,-7r_*/8]$, and the $Q_0$-good set lies inside $[-r_*/128,r_*/128]$.

If $\Delta\ge\mu R$, the minimizers of $Q_+$ and $Q_-$ are $R$ and $-R$.  For every $x\le R/2$,
\[
  Q_+(x)-Q_+(R)
  \ge Q_+(R/2)-Q_+(R)
  =\frac{\Delta R}{2}-\frac{3\mu R^2}{8}
  \ge\frac14S_\Delta.
\]
Thus the $Q_+$-good set lies in $(R/2,R]$, and symmetrically the $Q_-$-good set lies in $[-R,-R/2)$.  Meanwhile, for $|x|\ge R/2$,
\[
  Q_0(x)-Q_0(0)\ge\frac{\Delta R}{2}+\frac{\mu R^2}{8}>\frac14S_\Delta,
\]
so the $Q_0$-good set lies in $(-R/2,R/2)$.  The sets are pairwise disjoint in both regimes.

Under the uniform prior on the three functions, any fixed zero-query output belongs to at most one good set, so its average success probability is at most $1/3$.  Integrating over a randomized output and then fixing an objective proves the claim.
\end{proof}

\section{Calibration and completion of the proof}\label{sec:calibration}

If $\eps\le c_0S_\Delta$ with $c_0\le1/2$, then $m_\eps\ge1$, and $m_\eps\ge2$ whenever $d\ge2$.

We first record the elementary random-tape transfer used in both source branches.  Suppose a fixed prior $\cD$ has the property that every deterministic $T$-query algorithm has prior-averaged expected gap at least $a$.  Fixing each random tape and integrating by Tonelli shows that every randomized algorithm has a fixed objective in the support of $\cD$ with expected gap at least $a$.  Likewise, if every deterministic section has prior failure probability at least $p$, then every randomized algorithm has a fixed objective with failure probability at least $p$.  The common-prior property in \cref{thm:source-facts} is exactly what makes this step valid.

\begin{theorem}[All-dimensional expected-error lower bound]\label{thm:expected-lower}
There are universal constants $c,c_0>0$ such that, for every $d\in\mathbb N$, $R,\mu>0$, $L>\mu R$, and
\[
   0<\eps\le c_0S_\Delta,
\]
one has
\begin{equation}\label{eq:expected-main-lb}
  T^{\mathrm{exp}}_\eps(\cF^{\mathrm{sc}}_{L,\mu,R})
  \ge c\frac{d m_\eps}{\log(e m_\eps)}.
\end{equation}
\end{theorem}

\begin{proof}
Let $b>0$ be a common lower-bound constant for
\cref{prop:expected-transfer,prop:prob-transfer}. Choose
$0<\kappa\le\min\{1,c_{\mathrm{dim}},b^2/16,b/4\}$,
and then fix a universal threshold
$M_0\ge\max\{d_{\mathrm{src}},2/\kappa,2\}$.
This choice makes both transferred error scales at least $4\eps$
in the large-parameter case below. Split into two cases.

\emph{Large effective parameter: $m_\eps\ge M_0$.}
Set $k=\lfloor\kappa m_\eps\rfloor$.  Since $d\ge m_\eps\ge M_0$, the source dimension threshold is met, and
\[
  1\le k\le c_{\mathrm{dim}}d,
  \qquad
  \kappa m_\eps/2\le k\le\kappa m_\eps,
  \qquad
  \frac{dk}{\log(ek)}\gtrsim\frac{dm_\eps}{\log(em_\eps)}.
\]
Moreover, by the definition of $m_\eps$,
\[
  \frac{\Delta R}{\sqrt{k}}\ge\frac{\eps}{\sqrt\kappa},
  \qquad
  \frac{\Delta^2}{\mu k}\ge\frac{\eps}{\kappa}.
\]
The choice of $\kappa$ makes the expected gap in \cref{prop:expected-transfer} larger than $2\eps$.  Hence every deterministic algorithm with
\[
  T\le c\frac{dm_\eps}{\log(em_\eps)}
\]
has prior-averaged expected error larger than $2\eps$.  The random-tape transfer gives a fixed target objective with randomized expected error larger than $\eps$.

\emph{Bounded effective parameter: $1\le m_\eps<M_0$.}
Choose the universal constant $c$ small enough that
\[
  T\le c\frac{dm_\eps}{\log(em_\eps)}
  \quad\Longrightarrow\quad
  T\le\frac d4.
\]
Then \cref{prop:spherical-expected} gives a fixed spherical objective with expected gap at least $S_\Delta/2>\eps$, after reducing $c_0$ if necessary.  This includes $d=1$, where the admissible integer budget is zero for a sufficiently small universal $c$.
\end{proof}

\begin{theorem}[All-dimensional constant-success lower bound]\label{thm:prob-lower}
There are universal constants $c,c_0>0$ such that, for every $d\in\mathbb N$, $R,\mu>0$, $L>\mu R$, and
\[
   0<\eps\le c_0S_\Delta,
\]
one has
\begin{equation}\label{eq:prob-main-lb}
  T^{\mathrm{prob}}_\eps(\cF^{\mathrm{sc}}_{L,\mu,R})
  \ge c\frac{d m_\eps}{\log(e m_\eps)}.
\end{equation}
\end{theorem}

\begin{proof}
Use the same constants $\kappa,M_0$ and the same split, and take
$c_0\le1/128$.

If $m_\eps\ge M_0$, choose $k=\lfloor\kappa m_\eps\rfloor$.  The calibration above and \cref{prop:prob-transfer} imply that every deterministic algorithm with
\[
  T\le c\frac{dm_\eps}{\log(em_\eps)}
\]
has target error larger than $\eps$ with prior probability at least $5/8>1/2$.  Integrating over a randomized algorithm's independent tape yields a fixed objective with failure probability at least $5/8$.

Suppose $m_\eps<M_0$ and $d\ge2$.  Choose $c$ small enough that the displayed budget is at most $d/256$.  By \cref{prop:spherical-prob}, every deterministic section fails with probability at least $2/3$ under the spherical prior, at error threshold $3S_\Delta/8$.  Taking $c_0<3/8$ and integrating over the random tape gives a fixed objective with failure probability greater than $1/2$.

Finally, let $d=1$.  Then $m_\eps=1$.  Choose $c<1$, so the largest integer query budget not exceeding $c d m_\eps/\log(em_\eps)=c$ is zero.  The three-function prior in \cref{lem:one-dimensional} makes every randomized zero-query algorithm fail on a fixed objective with probability at least $2/3$.  This proves the claimed real-valued lower bound because the successful query count is integer.
\end{proof}


\begin{proof}[Proof of \cref{thm:main}]
Necessity of $L\ge\mu R$ and the boundary rigidity follow from
\cref{thm:slack-geometry}. Conversely, $F(x)=\mu\norm{x}^2/2$
is admissible whenever $L\ge\mu R$, proving nonemptiness.
The zero-query statements follow from \cref{cor:zero-query}.
Choose $c_0$ small enough for both lower-bound theorems above and
at most one. Combining them with \cref{thm:upper} proves
\eqref{eq:main-two-sided}. If an accuracy term attains $m_\eps$,
the corresponding upper branch is $O(d m_\eps)$; if $m_\eps=d$,
the value-only branch is $O(d m_\eps\logfac)$. Since $\logfac>1$
in the stated range, the final envelope follows.
\end{proof}

\begin{corollary}[A separately specified residual bound]\label{cor:residual}
For $G,\mu,R>0$, let
\[
 \cF^{\mathrm{res}}_{G,\mu,R}
 =\left\{h+\frac{\mu}{2}\norm{\cdot}^2:
        h\text{ is convex and }G\text{-Lipschitz on }\B_R^d\right\}.
\]
Define $S_G$ and $m_\eps(G)$ by replacing $\Delta$ with $G$ in
\eqref{eq:S-delta} and \eqref{eq:m-eps}. For every $d\ge1$ and
$0<\eps\le c_0S_G$, both criteria satisfy
\[
 c\frac{d m_\eps(G)}{\log(e m_\eps(G))}
 \le T_\eps^{\mathsf{risk}}(\cF^{\mathrm{res}}_{G,\mu,R})
 \le C d m_\eps(G)\log(d+1)\log\!\frac{128S_G}{\eps}.
\]
The zero-query minimax expected error is $S_G$.
\end{corollary}

\begin{proof}
The residual class is contained in
$\cF^{\mathrm{sc}}_{G+\mu R,\mu,R}$, whose slack is $G$.
This gives the upper bound and the zero-query upper guarantee.
The regularized support functions in \eqref{eq:hard-family}, the
spherical family in \eqref{eq:spherical-family}, and the
one-dimensional packing all have residual Lipschitz constant at
most $G$ when $\Delta=G$. They therefore establish the same lower
bounds within the residual class. The zero-query spherical argument
also gives its exact expected-risk threshold.
\end{proof}

\section{Discussion}\label{sec:discussion}

The total Lipschitz constraint determines both the scale of the
unknown residual and the part of the domain in which optimization
needs to take place. This explains why $L-\mu R$, rather than $L$
alone, controls a theorem that remains valid near the compatibility
boundary. The same localization improves the accuracy logarithm in
the dimension branch. At a fixed relative accuracy $\eps/S_\Delta$,
the upper bound has no additional logarithmic deterioration as
$L\downarrow\mu R$.

The lower-bound reduction separates the information cost of exact
values from the geometry of strong convexity. The support-function
prior supplies the information obstruction. The radial profile
$\Psi$ supplies the conversion to a strongly convex objective.
The spherical family has a different role: it gives a direct
linear-in-dimension obstruction when the effective parameter is
bounded. Its enlarged-measurement proof explicitly accommodates
redundant adaptive queries and therefore applies without a rank
assumption.

Several aspects remain outside the result. The logarithmic factors
in \eqref{eq:main-two-sided} do not match, especially at very high
accuracy. The success-$1/2$ criterion does not by itself describe
optimal dependence on an arbitrary confidence parameter. The
constant-width accuracy band between $c_0S_\Delta$ and $S_\Delta$
is also not assigned an exact query complexity here. Finally,
extensions to other domains require additional geometric information:
the radial and local residual estimates exploit the Euclidean ball.

Together, the results give a complete uniform description of the
polynomial complexity profile in the stated accuracy range, with
the external information-theoretic inputs and the remaining precision
limitations explicit.

\begin{thebibliography}{9}
\bibitem{DornEtAl2026}
Y.~Dorn, D.~Dvinskikh, T.~Loginov, O.~Osmanov, A.~Shestakov,
N.~Tupitsa, and A.~Gasnikov.
\newblock Lower bounds for gradient-free convex optimization and convex-concave
  saddle-point problems.
\newblock \emph{arXiv preprint arXiv:2610.00345v1}, 2026.
\newblock \url{https://arxiv.org/abs/2610.00345v1}.

\bibitem{Kerger2026}
P.~Kerger.
\newblock Closing the oracle-complexity gap in derivative-free convex
  optimization: A near-quadratic lower bound from exact function values.
\newblock \emph{arXiv preprint arXiv:2607.13335v1}, 2026.
\newblock \url{https://arxiv.org/abs/2607.13335v1}.

\bibitem{LeeSidfordVempala2018}
Y.~T. Lee, A.~Sidford, and S.~S. Vempala.
\newblock Efficient convex optimization with membership oracles.
\newblock In \emph{Proceedings of the 31st Conference on Learning Theory},
  volume~75 of \emph{Proceedings of Machine Learning Research},
  pages 1292--1294, 2018.
\newblock \url{https://proceedings.mlr.press/v75/lee18a.html}.

\bibitem{Protasov1996}
V.~Yu. Protasov.
\newblock Algorithms for approximate calculation of the minimum of a convex
  function from its values.
\newblock \emph{Mathematical Notes}, 59(1):69--74, 1996.
\newblock \href{https://doi.org/10.1007/BF02312467}{doi:10.1007/BF02312467}.

\bibitem{Shamir2017}
O.~Shamir.
\newblock An optimal algorithm for bandit and zero-order convex optimization
  with two-point feedback.
\newblock \emph{Journal of Machine Learning Research}, 18(52):1--11, 2017.
\newblock \url{https://jmlr.org/papers/v18/16-632.html}.

\bibitem{ZhangEtAl2026}
H.~Zhang, C.~Zhang, Z.~Qi, and Z.~Lin.
\newblock Near-optimal lower bounds for randomized algorithms in exact value
  zeroth-order convex optimization.
\newblock \emph{arXiv preprint arXiv:2607.16558v2}, 2026.
\newblock \url{https://arxiv.org/abs/2607.16558v2}.

\end{thebibliography}

\end{document}
